% ---------------------------------------------------------------------------
% Kapitel 1: Einleitung
% Wortgleich uebernommen aus der Confluence-Seite
% "Kapitel 1: Einleitung (Entwurf)" (Space BEH, Seiten-ID 34504705),
% Stand 04.10.2026 (die vorherige LaTeX-Datei trug noch Version 7 vom
% 02.09.2026). Einzige Aenderung: Der Verweis auf den Befundstand der
% Referenzplattform zeigt auf die zusammengefuehrte Tabelle 3.1 der
% Kurzfassung statt auf die fruehere Tabelle 3.2.
% ---------------------------------------------------------------------------

\chapter{Einleitung}
\label{kap:einleitung}

\section{Problemfeld}
\label{sec:problemfeld}

An der Hochschule Bremen entstehen in jedem Semester studentische Projektarbeiten in großer Zahl. Ihr Verbleib nach der Bewertung ist ungeregelt: Ein Teil liegt auf den privaten Laufwerken der Bearbeitenden, ein anderer in einem Lernmanagementsystem, das auf den Kursbetrieb ausgelegt ist und nicht auf die Darstellung nach außen. Beide Ablagen dienen der Benotung und sind nicht darauf ausgelegt, die Arbeiten danach zu zeigen. Ein Ort, an dem abgeschlossene Projekte dauerhaft und öffentlich einsehbar wären, besteht nicht. Diese Ausgangslage beruht auf einer Rücksprache mit der Fakultät 4 (Elektrotechnik und Informatik), in der auch Zielgruppe und Anforderungen des Portals besprochen wurden (J. Egharevba, persönliche Mitteilung, Mai 2026).

Aus dieser Lage folgt ein Verlust auf mehreren Seiten. Studierenden fehlt ein Beleg der eigenen Arbeit, auf den sich in einer Bewerbung verweisen ließe. Der Hochschule entgeht eine Außendarstellung, die aus der Lehre selbst hervorginge und nicht eigens hergestellt werden müsste. Arbeitgebende und die fachlich interessierte Öffentlichkeit haben keinen niedrigschwelligen Zugang zu dem, was in den Studiengängen bearbeitet wird. Studierenden späterer Jahrgänge fehlt der Blick auf das, was vor ihnen entstanden ist, und damit ein Maßstab für die eigene Arbeit.

Lorenzo und Ittelson (2005) beschreiben ein E-Portfolio als digitalisierte Sammlung von Artefakten, die eine Einzelperson, eine Gruppe oder eine Institution repräsentiert, und sie unterscheiden studentische, lehrbezogene und institutionelle Portfolios; einschlägig ist die dritte Kategorie, denn Träger wäre die Hochschule selbst. Aus dem Anspruch, etwas zu repräsentieren, folgt für das vorliegende Vorhaben eine erste Festlegung: Eine geordnete Dateiablage genügt nicht. Ein Bestand, der die Einrichtung vertreten soll, muss gepflegt werden, bevor er dargestellt werden kann, und beides gehört in dasselbe System.

Dieselben Autoren beschreiben das institutionelle Portfolio zugleich als Zusammenführung studentischer und lehrbezogener Portfolios; ein Portal, das studentische Arbeiten sammelt und nach außen zeigt, ist danach beides. Es unterscheidet zwei Gegenstände: das einzelne Projekt, das eine Prüfung durchläuft und danach öffentlich wird, und die Profilseite einer Person, auf der die veröffentlichten Projekte zusammenlaufen und die über die Studienzeit hinweg wächst. Diese Trennung leitet Abschnitt~\ref{sec:eportfolio} aus einer eigenen Ordnung der Komponenten nach Mahasneh (2020) her.

Zur Sicht der Studierenden liegen empirische Daten vor. Mahasneh (2020) entwickelt aus einer Literaturdurchsicht und Interviews mit 20 Fachleuten ein Modell aus zehn Komponenten, das von der Biografie über Berichte und Projekte bis zu wissenschaftlichem Material und Proben studentischer Leistung reicht. 90 Studierende, die damit gearbeitet haben, beurteilen es mit einem Mittelwert von 4,58 auf einer fünfstufigen Skala positiv. Welche Bestandteile ein hochschulisches Portfolio enthalten sollte, lässt sich demnach aus vorhandener Forschung ableiten.

Eine positive Einstellung bedeutet noch nicht, dass ein System auch benutzt wird. Davis (1989) führt im Technologieakzeptanzmodell (TAM) die Akzeptanz einer Technologie auf zwei Wahrnehmungen zurück: die wahrgenommene Nützlichkeit und die wahrgenommene Benutzerfreundlichkeit. Shroff et al. (2011) haben das Modell auf E-Portfolio-Systeme übertragen. Für ein Portal, dessen Inhalte freiwillig eingestellt und vor der Veröffentlichung geprüft werden, ist Akzeptanz deshalb keine nachgelagerte Frage der Oberfläche. Bleibt das Einreichen mühsam oder sein Nutzen unklar, füllt sich der Bestand nicht, und die Sichtbarkeitslücke besteht fort.

\section{Ziele der Arbeit}
\label{sec:ziele}

Ziel der vorliegenden Arbeit ist die Konzeption und die prototypische Realisierung eines institutionellen E-Portfolio-Portals für die Hochschule Bremen, das studentische Projekte strukturiert darstellt und sie über Suche und Filter auffindbar macht. Studierende sollen einzelne Projekte in einem blockbasierten Editor aufbauen und zur Prüfung einreichen; über die Veröffentlichung entscheiden Redaktion und Administration. Daneben sollen sie eine Profilseite pflegen, auf der die veröffentlichten Projekte erscheinen, die sie verfasst oder an denen sie mitgewirkt haben. Für die Öffentlichkeit entsteht ein Zugang ohne Anmeldung. Der redaktionelle Freigabeprozess richtet sich dabei nicht gegen die Studierenden, sondern hält die Verantwortung für veröffentlichte Inhalte bei der Hochschule.

Das Portal tritt nicht an die Stelle des Lernmanagementsystems und bildet weder Kursbetrieb noch Bewertung ab; Gegenstand ist die Darstellung abgeschlossener Arbeiten. Ebenso wenig zielt die Untersuchung auf einen Marktvergleich verfügbarer E-Portfolio-Systeme, auf die Einbindung in die Verwaltungsinfrastruktur der Hochschule oder auf die organisatorische Einführung im Hochschulbetrieb. Der Bearbeitungszeitraum von neun Wochen begrenzt zusätzlich den Funktionsumfang: Umgesetzt wird, was zur Beantwortung der Teilfragen erforderlich ist, während weitergehende Funktionen dem Ausblick (Abschnitt~\ref{sec:ausblick}) vorbehalten bleiben.

Die übergeordnete Forschungsfrage dieser Arbeit lautet:

Welche funktionalen und gestalterischen Anforderungen muss ein hochschulisches E-Portfolio-Portal erfüllen, um studentische Projekte strukturiert zu präsentieren, auffindbar zu machen und von den beteiligten Nutzergruppen akzeptiert zu werden?

Diese Frage wird durch drei Teilfragen (TF) operationalisiert:

\begin{itemize}[topsep=0.4em, itemsep=0.4em]
	\item TF 1: Welche Anforderungen an Struktur, Inhalt und Workflow lassen sich aus der bestehenden Forschungsliteratur zu E-Portfolios ableiten?
	\item TF 2: Welche funktionalen und gestalterischen Merkmale fördern die Auffindbarkeit und Benutzerfreundlichkeit eines institutionellen Portfolio-Portals?
	\item TF 3: Wie lässt sich ein blockbasiertes Inhaltsmodell mit integriertem Redaktionsworkflow realisieren, das Studierenden eine flexible und medienreiche Projektpräsentation ermöglicht?
\end{itemize}

Beantwortet werden die Teilfragen auf vier Wegen: durch die Auswertung der Forschungsliteratur, die heuristische Analyse einer bestehenden Referenzplattform, die Entwicklung des Systems und eine eigene empirische Erhebung am gebauten System. Die Entwicklungsarbeit ist methodisches Mittel zur Überprüfung dessen, was aus der Literatur gewonnen wurde, und kein Selbstzweck; die Erhebung prüft anschließend, ob das Gebaute die vorab gesetzten Kriterien einlöst. TF 1 richtet sich auf vorhandenes Wissen und beugt der Gefahr vor, Anforderungen aus dem eigenen Vorverständnis zu schöpfen; wo jede Teilfrage bearbeitet und beantwortet wird, zeigt Abschnitt~\ref{sec:aufbau}.

Damit die Antworten überprüfbar bleiben, sind die Maßstäbe vorab festgelegt und nicht erst aus dem Ergebnis gewonnen. Die Akzeptanz misst sich an den beiden Konstrukten des TAM, der wahrgenommenen Nützlichkeit und der wahrgenommenen Benutzerfreundlichkeit (Davis, 1989). Erhoben werden sie mit den zwölf Items der beiden zugehörigen Skalen, ergänzt um die zwei Items zur Nutzungsabsicht, mit denen Davis (1989) in seiner zweiten Untersuchung die künftige Nutzung erfragt. Da Davis für diese Skalen keine Normwerte berichtet, dient als Orientierungspunkt der Mittelpunkt der siebenstufigen Antwortskala, also 4,00: Werte darüber gelten als eher zustimmend, Werte darunter als eher ablehnend, ein Wert von genau 4,00 als unentschieden, und zwar für alle drei Skalen. Dieser Punkt ist eine Festlegung der vorliegenden Arbeit und kein Literaturwert (Abschnitt~\ref{sec:testkonzept}). Den inhaltlichen Zuschnitt prüft das Komponentenmodell von Mahasneh (2020). Für die Gebrauchstauglichkeit stehen zwei Maßstäbe nebeneinander: das Urteil Begutachtender anhand der zehn Usability-Heuristiken (Nielsen, 2024) und der Summenwert der System Usability Scale (SUS) nach Brooke (1996), für den ein Zielwert von 68 Punkten gilt. Dieser Wert liegt auf der Notenskala aus 241 industriellen Untersuchungen (Sauro \& Lewis, 2016, zitiert nach Lewis \& Sauro, 2018) in der Mitte der Note C und entspricht damit dem Median der Studienmittelwerte. Die Festlegung vorab begrenzt den Spielraum, einen Maßstab nachträglich dem Ergebnis anzupassen; wo eine Anwendungsregel erst nach der Messung nötig wird, weist Abschnitt~\ref{sec:bewertungsrahmen} sie als solche aus.

\section{Lösungsansatz}
\label{sec:loesungsansatz}

Der Weg zur Beantwortung beginnt bei der Literatur und endet bei der Messung am gebauten System. Ein Teil der Anforderungen an Struktur und Freigabe der Inhalte leitet sich aus der Forschungsliteratur ab: Lorenzo und Ittelson (2005) liefern die institutionelle Perspektive, Mahasneh (2020) den Blick auf die Bestandteile; die übrigen stammen aus dem Anwendungsfall. Den zweiten Weg bildet die heuristische Analyse einer bestehenden Referenzplattform, des öffentlich zugänglichen Portfolio-Portals der Hochschule für Gestaltung Schwäbisch Gmünd, nach dem Verfahren von Nielsen und Molich (1990) und mit der heute kanonischen Liste der zehn Heuristiken (Nielsen, 2024) als Raster.

Untersucht wurde allein die Sicht ohne Anmeldung, weil nur der öffentliche Bereich zugänglich war. Nicht erfüllt ist allein die siebte Heuristik, Flexibilität und Effizienz der Nutzung: Ein Suchfeld fehlt ganz, wer Studiengang und Schlagwort zugleich einschränken möchte, findet dafür keinen Weg, und nach dem Semester lässt sich gar nicht auswählen. Teilweise erfüllt ist unter anderem die Sichtbarkeit des Systemzustands, weil eine Brotkrumennavigation fehlt und allein Seitenkopf und Bedienelement zeigen, welche Einschränkung wirkt; den vollständigen Befundstand führt Tabelle~\ref{tab:heuristiken} in Abschnitt~\ref{sec:verwandte-arbeiten}.

Aus diesen Befunden entsteht eine eigene Anforderung an die Auffindbarkeit, die über die Referenzplattform hinausgeht. Sie umfasst ein Suchfeld über Titel, Kurzbeschreibung, Personennamen und Schlagworte, kombinierbare Facettenfilter, eine sichtbare Anzeige der aktiven Auswahl und die Möglichkeit, eine Einschränkung einzeln oder die ganze Auswahl in einem Schritt zurückzunehmen. Übernommen werden dagegen das einheitliche Aufbauschema der Projektseiten und die Erschließung über Studiengang und Person; das Semester, das die Referenzplattform nur am Projekt nennt, wird im eigenen System zur auswählbaren Achse.

Bevor gebaut wird, entsteht der Entwurf zunächst als Handskizzen und Wireframes in Miro, anschließend als Prototyp in Figma, sodass Gestaltungsentscheidungen vor der ersten Codezeile verhandelbar werden. In der Qualitätssicherung der Realisierung kommen zusätzlich sprachmodellgestützte Prüfschritte zum Einsatz (Abschnitt~\ref{sec:qualitaetssicherung}). Das Portal unterscheidet vier Akteurgruppen: Studierende, Redaktion, Administration und externe Nutzende. Das unterschriebene Exposé nennt die zweite Gruppe Betreuende beziehungsweise Lehrende; die vorliegende Arbeit benennt sie nach ihrer Aufgabe im Portal und nicht nach einer Stellung im Lehrbetrieb. Drei Gruppen sind an ein Konto gebunden; die externen Nutzenden haben keines und sind zugleich die Hauptzielgruppe. Diese Gewichtung leitet die vorliegende Arbeit aus dem äußeren der beiden Zwecke ab, die Lorenzo und Ittelson (2005) dem institutionellen Portfolio zuschreiben (Abschnitt~\ref{sec:eportfolio}); sie deckt sich mit der Einschätzung der Fakultät 4 (Abschnitt~\ref{sec:methodisches-vorgehen}).

Die Projektinhalte entstehen in einem blockbasierten Editor, der eine Seite in einzeln anordenbare Bausteine zerlegt, darunter Text, Bilder, Galerien, Video und Audio. Der redaktionelle Ablauf kennt fünf Zustände. Der reguläre Weg führt vom Entwurf über die Einreichung und die Genehmigung zur Veröffentlichung; den fünften Zustand, die Überarbeitung, erreicht ein Projekt nur, wenn Redaktion oder Administration es zurückgeben. Bearbeiten lässt sich ein Projekt allein im Entwurf und in der Überarbeitung, damit ein geprüfter Stand nicht unbemerkt verändert wird.

Das fertige System wird von den Verfassern inspiziert und mit elf Personen erprobt, neun Studierenden und zwei Externen; auch die Aufgaben der Redaktion und der Administration bearbeiten Studierende. Die Erhebung verbindet vier Methoden, von denen jede eine Frage beantwortet, welche die übrigen offenlassen. Beim lauten Denken, für das Nielsen (1994b) die methodische Referenz ist, bearbeiten die Teilnehmenden Aufgaben entlang der Kernpfade des Portals und zeigen dabei Bedienprobleme und ihre Ursachen. Die heuristische Evaluation, von den beiden Verfassern vor dem Nutzertest getrennt vorgenommen, erfasst auch Bereiche, zu denen keine Aufgabe führt. Die SUS misst die wahrgenommene Gebrauchstauglichkeit (Brooke, 1996), der Fragebogen zum TAM die wahrgenommene Nützlichkeit, die wahrgenommene Benutzerfreundlichkeit und die Nutzungsabsicht; Shroff et al. (2011) haben das Modell für E-Portfolio-Systeme mit einem eigenen, aus früheren TAM-Skalen abgeleiteten Fragebogen geprüft.

Lautes Denken, Heuristiken und SUS liegen nach der Annahme der vorliegenden Arbeit auf der Seite des Aufwands: Sie zeigen vor allem, wie mühsam die Benutzung ist. Ob ihr Ergebnis den Beteiligten etwas wert ist, zeigt erst die Skala zur wahrgenommenen Nützlichkeit, denn ein Portal kann leicht zu bedienen und zugleich entbehrlich sein. Beide Fragebögen werden deshalb getrennt ausgewertet und nicht ineinander umgerechnet.

\section{Aufbau der Arbeit}
\label{sec:aufbau}

Kapitel~\ref{kap:anforderungsanalyse} bestimmt Akteurgruppen und Anforderungen mit ihrer Herkunft; es greift dabei auf Quellen vor, die Kapitel~\ref{kap:grundlagen} im Einzelnen auswertet. Kapitel~\ref{kap:grundlagen} legt die Grundlagen aus der E-Portfolio-, der Akzeptanz- und der Usability-Literatur und berichtet die Analyse der Referenzplattform. Beide Kapitel bearbeiten TF 1, Abschnitt~\ref{sec:usability} liefert zudem die Kriterien für TF 2. Kapitel~\ref{kap:konzeption} entwirft das System; die Abschnitte~\ref{sec:inhaltsmodell} und~\ref{sec:workflow} tragen TF 3, Abschnitt~\ref{sec:ui-design} den gestalterischen Teil von TF 2. Kapitel~\ref{kap:realisierung} erprobt den Entwurf an der Realisierung. Kapitel~\ref{kap:test} berichtet Methodik und Daten des Tests ohne Wertung; Kapitel~\ref{kap:evaluation} bewertet sie, beantwortet in Abschnitt~\ref{sec:bewertung-test} TF 2 und gleicht beide Anforderungskataloge mit dem gebauten Stand ab. Abschnitt~\ref{sec:gesamtzusammenfassung} beantwortet TF 1 und TF 3, verdichtet die Antwort auf TF 2, führt die Forschungsfrage zusammen und benennt die Grenzen der Arbeit; Abschnitt~\ref{sec:ausblick} gibt den Ausblick.
